\section{単一$A_1$根壁の行列因子化と安定化したExt積}\label{sec:A1-matrix}
本節は前節までの標準弱楕円gluingのうち，反標準模型
$\overline S$がちょうど一つの$A_1$特異点$x$を持つ場合を扱う．
特異点を通らない滑らかな楕円境界$E$と$q\in\Pic^0(E)$は固定する．
\emph{ここでは$\Db(\qgr R_q)$の局所環を構成するのではなく}，
全$q$に共通な$\overline{\cB}\subset\overline{\cD}_q$の対象を
古典的な特異点で完備化して計算する．

\subsection{一周期的なrank-one加群}
\begin{lemma}[明示的行列因子化]\label{lem:A1-periodic-module}
$k=\mathbb C$，
\[
 A=k[[u,v,w]]/(uv-w^2)
  \simeq k[[s,t]]^{\{\pm1\}},
 \qquad u=s^2,\quad v=t^2,\quad w=st
\]
と置く．符号反転$s,t\mapsto-s,-t$に対する奇部分
$N=k[[s,t]]_{\rm odd}=As+At$はrank-oneの非自由MCM $A$加群で，
\[
 B=\begin{pmatrix}-w&-v\\ u&w\end{pmatrix},
 \qquad B^2=(w^2-uv)I_2
\]
に対して$N\simeq\operatorname{coker}(B:A^2\to A^2)$である．
従って
\[
 \cdots\xrightarrow{B}A^2\xrightarrow{B}A^2
 \xrightarrow{B}A^2\longrightarrow N\longrightarrow0
\]
は最小自由解消で，$\Omega_A N\simeq N$となる．
さらに
\[
 \underline{\End}_A(N)
 :=\End_A(N)/\{\text{自由加群を経由する写像}\}\simeq k.
\]
\end{lemma}
\begin{proof}
$k[[s,t]]$の偶部分と奇部分の分解から$N$は有限生成で，
$N$は正則環$k[[s,t]]$の$A$-直和因子なので深さ2のMCM加群である．
$A^2\to N$，$(a,b)\mapsto as+bt$の関係は
$(-w,u)$と$(-v,w)$で生成される．
実際$as+bt=0$なら，正則局所環$k[[s,t]]$上で
$(a,b)=(th,-sh)$，$h$は奇部分に属し，
$h=s\alpha+t\beta$（$\alpha,\beta\in A$）と書ける．
これが表示行列$B$を与える．
$B^2=(w^2-uv)I_2$は直接の行列積であり，
hypersurface上のmatrix factorizationの完全性から
表示した一周期的自由解消を得る\cite{EisenbudMF}．

$N$に$s$を掛けると$A$のfractional ideal
$I=(u,w)$に同型となる．$A$は正規なので，
rank-one反射的idealの自己準同型環は$(I:I)=A$である．
また
$I^{-1}=(u,w)/u$なので，rank-one加群$N$から自由加群を経由する写像の
$\End_A(N)=A$における像はtrace ideal
\[
 II^{-1}=(u,w)\,(u,w)/u=(u,v,w)=\mathfrak m_A.
\]
従って安定自己準同型環は$A/\mathfrak m_A=k$となる．
\end{proof}

\begin{proposition}[剰余体の第二syzygyと安定自己準同型代数]
\label{prop:A1-stable-end}
$k=A/\mathfrak m_A$に対して
\[
 \boxed{\Omega_A^2(k)\simeq N^{\oplus2}}
\]
が成り立つ．したがって$A$の特異性圏では
$\mathfrak m_A\simeq N^{\oplus2}$（非標準的な同型）があり，
$\eta\in\Hom_{D_{\rm sg}(A)}(N,N[1])$をshiftとの可逆同型から選ぶと
\[
\boxed{
 \bigoplus_{n\in\mathbb Z}
 \Hom_{D_{\rm sg}(A)}(\mathfrak m_A,\mathfrak m_A[n])
 \ \simeq\ M_2\bigl(k[\eta,\eta^{-1}]\bigr),
 \qquad\deg\eta=1
}
\]
が次数付き代数として成り立つ（非標準的同型）．
特に安定化した自己Extは全整数次数で4次元であり，
積を含めた一周期性を持つ．
\end{proposition}
\begin{proof}
$\Omega_A^2(k)$は$A^3\xrightarrow{(u,v,w)}A$のkernelである．
$S=k[[s,t]]$上で$(s^2,t^2,st)$のHilbert--Burch表示を
\[
 0\longrightarrow S^2
 \xrightarrow{
 \begin{pmatrix}-t&0\\0&-s\\s&t\end{pmatrix}}
 S^3\xrightarrow{(s^2,t^2,st)} (s,t)^2
 \longrightarrow0
\]
と書く．各列がsyzygyであり，$(s,t)^2$は高さ2の三生成idealなので，
この二列が全syzygyを生成する（Hilbert--Burchの初等的な場合）．
符号反転群の不変部分を取る操作は標数0でexactである．
上の二列は奇なので，$S^3$中の不変syzygyを得るには
係数が奇である必要がある．従って
\[
\ker(A^3\xrightarrow{(u,v,w)}A)
 \simeq S_{\rm odd}^{\oplus2}=N^{\oplus2}.
\]

$D_{\rm sg}(A)=\Db(\operatorname{mod}_{\rm fg}A)/\Perf(A)$では
自由加群が零なので，短完全列によるtriangleから
$\Omega^2 k\simeq k[-2]$，$\mathfrak m_A=\Omega k\simeq k[-1]$である．
補題\ref{lem:A1-periodic-module}より
$\Omega N\simeq N$，すなわち$N[1]\simeq N$なので，
$\mathfrak m_A\simeq N^{\oplus2}$となる．
Gorenstein局所環上のBuchweitz--Orlovの同値により，
$N$の特異性圏内の次数0自己準同型は
前補題の安定自己準同型環$k$に一致する
\cite{BuchweitzMCM,OrlovSingularity}．
$N[1]\simeq N$の可逆射$\eta$を一つ選べば，
全整数次数の自己準同型は$\eta^n$の一つのスカラー倍であり，
合成はLaurent多項式の積になる．二つのコピーを取って主張を得る．
\end{proof}

\begin{corollary}[$q$に依存しない局所積模型]\label{cor:A1-local-q-invariant}
$\overline{\cB}\subset\overline{\cD}_q$の特異点ideal sheaf
$I_x$の像$J_{x,q}$を，局所完備化して
$D_{\rm sg}(\widehat{\mathcal O}_{\overline S,x})$へ送ると，
その像の次数付き安定自己準同型代数は
\[
 M_2\bigl(k[\eta,\eta^{-1}]\bigr),\qquad\deg\eta=1
\]
であり，$q$に依存しない．
\end{corollary}
\begin{proof}
$I_x$の完備化は$\mathfrak m_A$であり，
$\overline{\cB}$の完全忠実な包含は全$q$で共通である．
前命題を適用すればよい．
\end{proof}

\begin{remark}[ordinary Ext，安定化，非可換標的の区別]
定理\ref{thm:duval-high-ext}の
$\Ext^n_{\overline{\cD}_q}(J_{x,q},J_{x,q})$は$n\ge3$で4次元である．
本節のLaurent代数の計算は\emph{局所完備化後にperfect複体を零とした}
$D_{\rm sg}(A)$の自己準同型の\emph{積}に関する計算であり，
普通のExtの全低次数を含むYoneda代数がLaurent代数だという主張ではない．
また$q\ne\OO_E$の$D^b(\qgr R_q)$に対応する対象や積を
この計算だけから同定したことにはならない．
非自明$q$の自然関手の完全忠実性・本質的全射性は依然として未解決であり，
次は$A_1$局所模型の周期射$\eta$が標的$qgr$側にどのように実現されるかを調べる．
\end{remark}
